
% ============================================================================
\clearpage
\section*{v4.10.0-r13 Source-Only: Status Notes on the Neutral Qubit and the Second Family Structure}
\addcontentsline{toc}{section}{v4.10.0-r13 Source-Only: Status Notes on the Neutral Qubit and the Second Family Structure}
\label{sec:v4100-r13-neutral-qubit-family}
% ============================================================================

\paragraph{Purpose.}
After v4.10.0-r12 the OP3 work continued on the profile of the missing structure: a colour- and charge-neutral qubit on
which the generation cycle acts as an element of order three, with a reference state under the magic angle
\(\arccos(1/\sqrt3)\).  This layer records what the explorations of 13 and 14 November 2026 found, corrects one sentence of
r12.4(c), and states where the search stops.  All items except H103 are explorations without pre-registration; none is
compared with data.  The notes are \texttt{notes/\allowbreak research/\allowbreak v4\_10\_0\_r13/} and
\texttt{notes/\allowbreak research/\allowbreak op3\_*\_2026\_11\_13.*}, \texttt{op3\_*\_2026\_11\_14.*}.  Earlier layers are not rewritten.

\noindent\textbf{Source-only status.}
This v4.10.0-r13 layer is source-only.  It does not build a PDF, does not assign a new DOI,
and does not authorize a public release.  The v5.0.0 Version DOI
\href{https://doi.org/10.5281/zenodo.24152615}{10.5281/zenodo.24152615} and the Concept DOI
\href{https://doi.org/10.5281/zenodo.19210728}{10.5281/zenodo.19210728} are unchanged by it.

\subsection*{v4.10.0-r13 status notes}
\label{subsec:v4100-r13-status}

\paragraph{Theoremlet v4.10.0-r13.1 (a neutral qubit exists under an assumed charge axis; correction to r12.4(c)).}
r12.4(c) says that no colour- and charge-neutral native carrier was found and that the sedenion cycle \(\psi\) moves the
charge axis.  That sentence holds for the inherited charge \(Q_S\) with axis \(e_7\).  A backward search over all charges on
\(\mathbb C\otimes\mathbb S\) with \([Q,\psi]=0\), \([Q,\mathfrak{su}(3)_c]=0\), \(\bar Q=1-Q\) and spectrum
\(\{0,\tfrac13,\tfrac23,1\}\) gives, up to two signs,
\[
Q_*=P_0\bigl(\tfrac12-\tfrac i2 L_{e_8}\bigr)P_0+\tfrac23P_3+\tfrac13P_{\bar3},
\]
with the charge axis \(e_8\), the only \(\psi\)-fixed imaginary direction.  Its kernel
\(N=\operatorname{span}(1+ie_8,\;e_7-ie_{15})\) is a colour- and charge-neutral qubit, and \(\psi\) acts on it as
\(\operatorname{diag}(1,\omega)\).  The condition \(Q|_{\OO}=Q_8\) is given up: \(Q_*\) does not leave \(\OO\) invariant.
The axis \(e_8\) is an assumption.  It is entered in the counterterm registry as placeholder CT-003 with three registry
guards; every statement below that uses \(Q_*\) or \(N\) inherits them.  H103a (pre-registration \texttt{45314b2}) is
refuted: the inherited charge \(Q_S\) acts on \(N\) as \(\tfrac12+\tfrac13\sigma_y\), perpendicular to the \(\psi\) axis, and fixes
no reference under the magic angle.

\paragraph{Theoremlet v4.10.0-r13.2 (Hurwitz units and the reference state; not blind).}
\(\mathbb H_n=\operatorname{span}(1,e_7,e_8,e_{15})\) is associative, and \(N\) is a right module over it: \(iR_{e_7}\),
\(iR_{e_{15}}\), \(iR_{e_8}\) act on \(N\) as a Pauli triple with \(e_8\) as the \(\psi\) axis.  The 24 Hurwitz units of
\(\mathbb H_n\) act from the right faithfully as the binary tetrahedral group \(2T\) (projectively \(A_4\)).  The \(\psi\) axis is
a twofold axis of this group, and all four threefold axes lie under \(|\cos|=1/\sqrt3\) to it.  An eigenstate of a Hurwitz
unit of order three, for instance of \((1+e_7+e_8+e_{15})/2\), has under \(\psi\) the orbit overlap \(\rho^2=\tfrac12\), hence
\(K=\tfrac23\), and Bargmann phase \(\pi/4\).  The target was known in advance; this is a construction and not a derivation of
the Koide relation.  \(\psi\) does not normalise \(2T\).

\paragraph{Theoremlet v4.10.0-r13.3 (generation content under \(Q_*\)).}
Under \((Q_*,\text{colour type},\psi)\) the space \(\mathbb C\otimes\mathbb S\) carries exactly two complete generations
(\(2\,\nu\), \(2\,e^+\), two triplets \(u\), two triplets \(\bar d\)).  The \(\psi\) sectors are not generations
(\(\psi=1\): \(\nu,e^+\); \(\omega\): \(\nu,u,\bar d\); \(\bar\omega\): \(e^+,u,\bar d\)).  \(\psi\) commutes with \(Q_*\) and colour
and therefore acts inside every generation as a phase; it is not a generation step.  \(N\) is the space of two neutrino
flavours.  The third generation is absent, as in the earlier sedenion analysis (\(\mathbb S=(1,1)\oplus(1,\mathrm{sgn})\oplus(7,2)\)
under \(G_2\times S_3\)).

\paragraph{Theoremlet v4.10.0-r13.4 (the right action of \(\mathbb H_n\) is a family symmetry).}
On all of \(\mathbb C\otimes\mathbb S\) the right multiplications \(R_{e_7}\), \(R_{e_8}\), \(R_{e_{15}}\) commute with
\(\mathfrak{su}(3)_c\) and with \(Q_*\) (numerically below \(3\cdot10^{-14}\)); \(L_{e_7}\) and \(L_{e_{15}}\) do not commute with
\(Q_*\).  On the neutral and on the coloured part the quaternion relations hold exactly, with opposite orientation:
\(R_xR_y=R_{yx}\) on the neutral part, \(R_xR_y=R_{xy}\) on the coloured part.  The Lie algebra generated is therefore
\(\mathfrak{su}(2)_{\rm leptons}\oplus\mathfrak{su}(2)_{\rm quarks}\), and the \(R_g\) of the Hurwitz units represent \(2T\) on the whole
space only with \(\bar g\) on the coloured part.  Every species (\(\nu\), \(e^+\), \(u\), \(\bar d\)) is a doublet (Casimir \(3/4\)).
\(\psi\) acts on this family algebra as the rotation of \(\operatorname{Im}\mathbb H_n\) about \(e_8\) by \(120^\circ\).
Limitation: for two generations a \(U(2)\) per species that commutes with colour and charge follows from Schur's lemma;
new is only its realisation by one native operation.  The generations form a doublet, not three axes, and the doublet is
carried by all species, not only by the Majorana block.

\paragraph{Theoremlet v4.10.0-r13.5 (the Majorana vector on \(N\) is not fixed).}
A symmetric bilinear form on \(N\) is \(M=(m\cdot\sigma)\varepsilon\) with \(m\in\mathbb C^3\), a vector in
\(\operatorname{Im}\mathbb H_n\otimes\mathbb C\); its two singular values are the two masses.  The native bilinear form of \(\mathbb S\)
vanishes on \(N\).  Words of length one and two in \(L_a\), \(R_a\) (\(a=e_7,e_8,e_{15}\)) give only real directions, along
\(e_8\) or perpendicular to it.  A real direction gives two equal masses, a null vector exactly one mass.  Invariance under
\(\psi\) admits only \(M_{11}\) (one massive, one massless neutrino); invariance under a Hurwitz unit of order three gives the
space diagonal with two equal masses.  With a third mass zero, \(K\) runs from \(1\) to \(\tfrac12\) and passes \(\tfrac23\) only
at a free mixing angle; that is an adjustment and no prediction.  Two flavours do not give three masses.  The pair channel
\(\nu^c\nu^c\) of v4.9.0-r82.4 is not evaluated here; a map between \(\mathbb C\otimes\mathbb S\) and
\(\mathbb C\otimes\mathbb H\otimes\OO\) is missing.  By v4.10.0-r01.3(d) that channel is repulsive and forms no condensate, so it
is not available as a source of this vector; the note of 13 November that calls it attractive
(\texttt{op3\_\allowbreak neutral\_\allowbreak su2\_\allowbreak octant\_\allowbreak 2026\_\allowbreak 11\_\allowbreak 13.md}, field search, item 2) used the superseded sign.

\paragraph{Theoremlet v4.10.0-r13.6 (third generation: the two routes are complementary).}
\begin{enumerate}[leftmargin=2em]
\item[(a)] Compare \((\mathbb S\oplus\OO,\ \psi\oplus1)\) with the triality carrier \((\OO^3,\tau)\), \(\tau(x,y,z)=(y,z,x)\).  The
multiplicities of the eigenvalues \((1,\omega,\bar\omega)\) agree on the coloured part, \((6,6,6)\); for \(u\) and \(\bar d\)
``doublet plus singlet'' is the same \(\mathbb Z_3\) representation as three slots.  On the neutral part they are \((4,1,1)\)
against \((2,2,2)\).  The difference is the plane \(\operatorname{span}(1,e_8)\), fixed by \(\psi\), which carries the charge axis and
the \(\psi\)-fixed neutrino state and has no counterpart in \(\OO^3\).
\item[(b)] The \(e_8\) charge cannot be made \(\tau\)-symmetric on three slots.  On the neutral part
\(\mathbb R^2\otimes\mathbb R^3\) a charge \(\tfrac12-\tfrac i2A\) needs \(A^2=-1\); with \(A=d(1\otimes J)+c(\text{cycle}\otimes X)\) and
the pair coupling \(X\) of \(L_{e_8}\), a reflection, the off-diagonal block of \(A^2\) is \(c^2\cdot1\neq0\).  \(Q_*\) is a pair
structure.
\item[(c)] With the slot charge \(Q_8\) the carrier \(\OO^3\) has a \(\tau\)-symmetric splitting \(2+1\) for all species: the
symmetric slot mode is the singlet, and on the doublet of the modes \((\omega,\bar\omega)\) the cycle is
\(\tau=\operatorname{diag}(\omega,\bar\omega)\in SU(2)\).  A real slot vector lies on the equator and gives \(\rho^2=\tfrac14\),
\(K=\tfrac12\).
\item[(d)] No operation at hand supplies a twofold axis under the magic angle on that doublet.  Monomial slot operations
(a permutation times phases per slot; these include \(\tau\), \(\rho\), the sign changes of \(A_4\) of v4.9.0-r55 and the
slot-wise action of \(\mathrm{Spin}(8)\) that commutes with colour) preserve \(2+1\) only as a number times a permutation; of
\(10\,368\) elements with twelfth roots of unity \(72\) preserve it, and on the doublet they rotate about the \(\tau\) axis or
about axes perpendicular to it.  \(A_4\) does not preserve \(2+1\) at all; under it the three slots are the irreducible
vector.  The required element superposes all three slots.  The triality maps themselves were not recomputed.
\end{enumerate}
Reason for the complementarity: the quaternion group needs a doublet made of two slots (route A, where
\(R_{e_7},R_{e_8},R_{e_{15}}\) are monomial and the Hurwitz unit is their half-integer sum), whereas the slot cycle needs
three slots and turns the doublet into the Fourier pair (route B).  With the present inventory the profile cannot be met.
Not examined: a new algebra structure on \(\OO^3\) in which sums of slot operations are operations.

\subsection*{v4.10.0-r13 status in the three frames}
\label{subsec:v4100-r13-frames}

\begin{itemize}[leftmargin=2em]
\item M and T: all statements are algebraic statements about carriers of the matter-frame readout; nothing is
transported to T.
\item P: not evaluated.
\end{itemize}

\subsection*{v4.10.0-r13 register entries}
\label{subsec:v4100-r13-register}

\begin{center}
\small
\begin{tabular}{@{}>{\raggedright\arraybackslash}p{0.34\textwidth}>{\raggedright\arraybackslash}p{0.58\textwidth}@{}}
\toprule
statement & status after v4.10.0-r13\\
\midrule
r12.4(c) ``no neutral native carrier'' & holds for \(Q_S\); under the assumed axis \(e_8\) the qubit \(N\) exists (r13.1)\\
charge axis \(e_8\) & assumption, placeholder CT-003, three registry guards (r13.1)\\
H103a & refuted (r13.1)\\
reference under the magic angle & eigenstate of a Hurwitz unit, a construction (r13.2)\\
\(\psi\) as generation step & no; a phase inside two generations (r13.3)\\
family symmetry of route A & right action of \(\mathbb H_n\), doublet, leptons and quarks separately (r13.4)\\
Majorana vector on \(N\) & not fixed; no three masses (r13.5)\\
profile of the missing structure & not met by the present inventory (r13.6)\\
H100 / H102 & pending / pending\\
\bottomrule
\end{tabular}
\end{center}

\subsection*{v4.10.0-r13 ledger and guard update}
\label{subsec:v4100-r13-ledger-guard}

\[
R_{\rm chk}^{v4100\text{-}r13}=(1,1,0)
\]
for (status notes recorded; formulas checked in the cited notes; derivation of \(K\), third generation and native reference
open).

The layer spends no guard:
\[
R_{\rm glob}^{(v4100\text{-}r13,+)}=(0,1,1,1,1),
\qquad
\Gamma_{\rm guard}^{v4100\text{-}r13}=(1,1,1,1,1),
\qquad
\beta_4=1/256\ \text{unspent}.
\]
OP~3 and OP~4 remain open.  The layer does not imply exact finite \(C5\) zero closure, exact all-readout terminality,
\(C6/T2\), a global photon-clock/RH theorem, or global all-readout closure.
